The Reock score was introduced by \citet{reock1961note} as a measure of how
much of the minimum enclosing circle a district fills.
Its key conceptual insight is that the minimum bounding circle (MBC) is determined
entirely by the most extreme points of a district---its ``tips''---rather than
by the full boundary perimeter.
This makes Reock insensitive to boundary digitisation artifacts, coastline
irregularities, and river boundaries that strongly affect perimeter-based
metrics such as Polsby-Popper.

\subsection{Reock in Court}

Reock has been cited in redistricting litigation in several states:
\begin{itemize}
  \item \textbf{North Carolina}: Expert reports in \textit{Harper v.\ Hall} (2022)
    cited Reock scores alongside PP values.
    The NC state supreme court's compactness analysis referenced the general
    principle that compact districts should ``fill'' their bounding geometry~\citep{harper2022}.
  \item \textbf{Wisconsin}: Expert reports in \textit{Gill v.\ Whitford} (2018)
    included Reock computations for Assembly districts under the challenged
    legislative plan~\citep{gill2018}.
  \item \textbf{Georgia}: Redistricting litigation in Georgia has produced
    expert reports using Reock for House and Senate districts.
  \item \textbf{Florida}: The Florida Supreme Court-ordered redistricting (2015--2016)
    included Reock-based analysis by court-appointed experts.
\end{itemize}

\subsection{Why Reock for Coastal States}

For states with Atlantic, Gulf, or Pacific coastlines, PP is unreliable because
coastal boundaries contribute zigzag perimeter without corresponding district
``mass.''
Reock avoids this problem: the MBC is determined by the extreme points
(typically inland or cross-bay points), and the jagged coastline between
those extreme points does not affect the MBC.
This paper provides the first systematic comparison of Reock boundary sensitivity
against PP boundary sensitivity under controlled computational conditions.

\subsection{Moving-Knife Algorithm}

The moving-knife algorithm (MKA) is a bisect structure algorithm designed to
maximise the minimum Reock score across all districts.
It operates by selecting bisection cuts that minimise the MBC radius growth,
preferring cuts that produce round sub-regions rather than elongated ones.
We provide the first empirical validation of MKA's Reock maximisation claim
on NC 2020 data.
